%! AP Statistics — Unit 1, Lesson 1.7 — Group Activity ANSWER KEY
\documentclass[10pt]{article}
\usepackage{apstats-article}
\usepackage{apstats-key}

%=============================================================================
\begin{document}

\pageheader{Unit 1, Lesson 1.7}{Group Activity --- Answer Key}

\vspace{0.10in}

\begin{tcolorbox}[colback=sky, colframe=navy, boxrule=0.8pt, arc=2mm,
    left=3mm, right=3mm, top=2mm, bottom=2mm]
\small
\textbf{Part 1} (group, 8 min): Calculate summary statistics for a swim team dataset.\\
\textbf{Part 2} (group, 7 min): Detect and discuss an outlier; choose appropriate statistics.\\
\textbf{Part 3} (individual, 5 min): Interpret a second dataset and make a recommendation.
\end{tcolorbox}

\vspace{0.10in}

% ── PART 1 ───────────────────────────────────────────────────────────────────
{\large\bfseries\color{navy} Part 1 \quad Calculate Summary Statistics}

\vspace{0.10in}

\noindent\small A swim team coach recorded the lap times (seconds) for 15 swimmers
during practice. The sorted data are:

\vspace{0.08in}
\begin{center}
\small
62,\enspace 64,\enspace 65,\enspace 67,\enspace 68,\enspace 69,\enspace 70,\enspace
71,\enspace 72,\enspace 73,\enspace 74,\enspace 75,\enspace 77,\enspace 80,\enspace 95
\end{center}

\vspace{0.12in}

\begin{enumerate}[leftmargin=1.5em, itemsep=8pt]

  \item Find the \textbf{mean}. Show the sum and division.\\[5pt]
        Sum $= \ans{1082}$ \qquad $\bar{x} = \ans{$1082/15 \approx 72.1$ s}$

  \item Find the \textbf{median}. Which position holds the median?\\[5pt]
        Position: \ans{8th} \qquad Median $= \ans{71}$ seconds

  \item Find $Q_1$ and $Q_3$ by finding the median of each half.\\[5pt]
        Lower half (7 values): \ans{62, 64, 65, 67, 68, 69, 70}\\[4pt]
        $Q_1 = \ans{67}$ \hfill
        Upper half (7 values): \ans{72, 73, 74, 75, 77, 80, 95}\\[4pt]
        $Q_3 = \ans{75}$ \hfill
        $\text{IQR} = Q_3 - Q_1 = \ans{8}$ seconds

  \item Find the \textbf{range}.\\[5pt]
        Range $= \ans{95} - \ans{62} = \ans{33}$ seconds

\end{enumerate}

\vspace{0.12in}

% ── PART 2 ───────────────────────────────────────────────────────────────────
{\large\bfseries\color{navy} Part 2 \quad Outlier Detection \& Choosing Statistics}

\vspace{0.10in}

\begin{enumerate}[leftmargin=1.5em, itemsep=8pt, resume]

  \item Calculate the lower and upper fences using the $1.5 \times \text{IQR}$ rule.\\[5pt]
        Lower fence $= Q_1 - 1.5(\text{IQR}) = \ans{$67-1.5(8)=55$}$\\[4pt]
        Upper fence $= Q_3 + 1.5(\text{IQR}) = \ans{$75+1.5(8)=87$}$

  \item Is the value 95 an outlier? Justify your answer.\\[5pt]
        \ansline{Yes. Since $95 > 87$ (the upper fence), 95 seconds is an}\\[1pt]
        \ansline{outlier by the $1.5\times\text{IQR}$ rule.}

  \item Compare the mean and median you calculated.  Which is larger?
        What does this tell you about the shape of the distribution?\\[5pt]
        \ansline{Mean $(72.1) >$ median $(71)$, so the distribution is slightly}\\[1pt]
        \ansline{\textbf{skewed right}; the outlier at 95 pulls the mean upward.}

  \item Based on the shape and the outlier, which pair of statistics should the
        coach use to describe the typical lap time and its variability?
        Circle your choice and write one sentence justifying it.

        \vspace{8pt}
        \hspace{2em}\textbf{Mean + Standard Deviation}
        \hspace{4em}\textbf{Median + IQR}
        \hspace{1em}\ans{$\leftarrow$ correct}

        \vspace{8pt}
        \ansline{The distribution is skewed right and contains an outlier (95),}\\[1pt]
        \ansline{so median (71 s) + IQR (8 s) are resistant and more accurate.}

\end{enumerate}

\vspace{0.12in}

% ── PART 3 ───────────────────────────────────────────────────────────────────
{\large\bfseries\color{navy} Part 3 \quad Individual Interpretation}

\vspace{0.10in}

\noindent\small A different coach reports: ``The mean number of freestyle laps
completed per week for my 20 swimmers is 42, and the median is 44. The IQR is 8
and the standard deviation is 12.''

\begin{enumerate}[leftmargin=1.5em, itemsep=8pt, resume]

  \item What does the fact that mean $<$ median suggest about the shape of this
        distribution? Explain.\\[5pt]
        \ansline{Mean $<$ median suggests the distribution is \textbf{skewed left} ---}\\[1pt]
        \ansline{a few unusually low values pull the mean below the median.}

  \item The coach wants to tell parents ``how many laps a typical swimmer
        completes per week.'' Which single number should the coach report,
        and why?\\[5pt]
        \ansline{The \textbf{median (44 laps)}, because the distribution is skewed}\\[1pt]
        \ansline{left and the median resists those low values.}

\end{enumerate}

\end{document}
